\section{Experiments}
\label{sec:experiments}

\subsection{Experimental Setup}
\label{sec:setup}

\paragraph{Benchmark and model.}
We evaluate on the LIBERO-Goal suite~\cite{liu2023_libero}, comprising 10 tabletop manipulation tasks with 50 human demonstrations each.
We use OpenVLA-7B~\cite{kim2024_openvla,karamcheti2024}, a vision-language-action model consisting of a SigLIP-SO400M vision encoder~\cite{zhai2023_siglip}, an MLP projector, and a Llama-2 7B language backbone, fine-tuned on the LIBERO-10 split.

\paragraph{Representation sources.}
We extract representations from five sources:
(a)~\textbf{SigLIP+Projector}---the vision encoder output after MLP projection into LLM input space;
(b)~\textbf{LLM image\_mean}---the mean of image token activations at each LLM layer, capturing distributed visual information;
(c)~\textbf{LLM last\_preaction}---the final token position before action decoding (readout token) at each LLM layer;
(d)~\textbf{DINOv2 ViT-L/14}~\cite{oquab2024_dinov2} \texttt{cls\_token}; and
(e)~\textbf{CLIP ViT-L/14}~\cite{radford2021_clip} \texttt{patch\_mean}.
Sources (d) and (e) serve as external baselines with no VLA training signal.

\paragraph{Experimental conditions.}
We use a $2 \times 2$ factorial design crossing \emph{training status} (trained vs.\ untrained) with \emph{instruction presence} (with vs.\ without language instruction).
The untrained control applies \texttt{from\_config()} to reinitialize the LLM parameters (6.74B random weights) while preserving the frozen SigLIP encoder (0.80B pretrained weights), isolating the LLM's learned contribution from the vision encoder's pretrained geometry.

\paragraph{Representational Similarity Analysis (RSA).}
Following~\cite{kriegeskorte2008}, we compute pairwise representational dissimilarity matrices (RDMs) for each representation source and for the 7-dimensional action space.
We report Spearman rank correlation $\rho$ between representation and action RDMs, with 95\% confidence intervals from 10{,}000 bootstrap resamples.

\paragraph{Functional probes.}
We train MLP probes (input $\to$ 256 $\to$ 7) with Adam (lr${}=10^{-3}$, batch size 256, 50 epochs, early stopping patience 5) to predict 7-D action vectors from frozen representations~\cite{alain2017}. We report mean $R^2$ across action dimensions.

%------------------------------------------------------------------
\subsection{Vision Encoder Dominance}
\label{sec:vision_dominance}

This experiment tests whether vision-language pretrained encoders provide the dominant source of action-relevant geometric structure in VLA representations.

Table~\ref{tab:rsa_overview} reports RSA $\rho$ between each representation source and the action space.
The SigLIP+Projector output achieves the highest alignment ($\rho = 0.615$, 95\% CI $[0.526, 0.695]$), followed by CLIP ViT-L/14 ($\rho = 0.557$, $[0.442, 0.660]$).
Both vision-language pretrained encoders substantially exceed DINOv2 ($\rho = 0.388$), a vision-only model.
The VLA's LLM layers (best image\_mean: $\rho = 0.501$ trained, $\rho = 0.485$ untrained) fall between VL-pretrained encoders and DINOv2 (Figure~\ref{fig:rsa_bar}).

Three findings emerge.
First, VL pretraining---not LLM processing---is the primary source of action-relevant structure.
SigLIP and CLIP, both pretrained on image-text pairs, achieve $\rho \geq 0.557$, while DINOv2, pretrained with self-supervised vision objectives alone, reaches only $\rho = 0.388$.
This 0.169-point gap indicates that language-grounded visual features carry substantially more action-relevant geometry than features learned without language, consistent with behavioral evidence of visual dominance in VLA policies~\cite{grant2026}.

Second, the MLP projector provides incremental alignment: SigLIP+Projector ($0.615$) exceeds CLIP ($0.557$) by $\Delta = 0.058$, consistent with the projector mapping vision features into the LLM's embedding space.

Third, the LLM degrades rather than enhances vision encoder geometry.
Processing through the LLM reduces RSA from $0.615$ (projector output) to $0.501$ (best LLM image\_mean), a drop of $\Delta = 0.114$.

\paragraph{Geometric-functional dissociation.}
Functional probes reveal a dissociation between geometric similarity and functional decodability.
DINOv2 achieves $R^2 = 0.712$, statistically indistinguishable from SigLIP+Projector ($R^2 = 0.706$, $p = 0.22$, n.s.), despite its substantially lower RSA ($0.388$ vs.\ $0.615$).
High linear decodability does not require high geometric alignment with the action space---DINOv2 features are functionally rich but geometrically dissimilar to actions.

\paragraph{Metric sensitivity.}
This dissociation extends to alternative similarity metrics.
Linear CKA~\cite{kornblith2019} produces a different ranking from RSA: untrained LLM ($\text{CKA} = 0.089$) $>$ trained LLM ($0.083$) $>$ SigLIP+Projector ($0.081$) $>$ DINOv2 ($0.066$), partially reversing the RSA ordering (Figure~\ref{fig:metric_comparison}).
RDM kernel CKA restores the RSA-consistent ranking (Proj $>$ Trained $>$ Untrained $>$ DINOv2).
This divergence reflects the known sensitivity of representation similarity metrics to different geometric properties and underscores the importance of multi-metric evaluation in representation analysis.

\begin{table}[t]
\centering
\caption{RSA between representation and action RDMs ($\rho$ $\uparrow$). VL-pretrained encoders (SigLIP, CLIP) dominate over vision-only DINOv2 and LLM-processed representations. 95\% CIs from 10{,}000 bootstrap resamples.}
\label{tab:rsa_overview}
\small
\begin{tabular}{llcc}
\toprule
Source & Aggregation & RSA $\rho$ & 95\% CI \\
\midrule
SigLIP+Projector & --- & 0.615 & {[0.526, 0.695]} \\
CLIP ViT-L/14 & patch\_mean & 0.557 & {[0.442, 0.660]} \\
\midrule
VLA Trained w/ inst. & image\_mean & 0.501 & --- \\
VLA Trained no inst. & image\_mean & 0.501 & --- \\
VLA Untrained w/ inst. & image\_mean & 0.485 & --- \\
VLA Untrained no inst. & image\_mean & 0.482 & --- \\
\midrule
DINOv2 ViT-L/14 & cls\_token & 0.388 & --- \\
\bottomrule
\end{tabular}
\end{table}

\begin{figure}[t]
\centering
\includegraphics[width=\columnwidth]{figures/fig2.pdf}
\caption{RSA with the action space ($\rho$) across representation sources. VL-pretrained encoders (SigLIP $\rho{=}0.615$, CLIP $\rho{=}0.557$) substantially exceed both LLM-processed representations ($\rho \leq 0.501$) and vision-only DINOv2 ($\rho{=}0.388$). Hatched bars indicate no-instruction conditions. Error bars: 95\% bootstrap CIs. Dashed line: DINOv2 reference.}
\label{fig:rsa_bar}
\end{figure}

%------------------------------------------------------------------
\subsection{LLM Redistributes, Does Not Enhance, Action Structure}
\label{sec:llm_redistribution}

This experiment tests how the LLM transforms the action-relevant structure inherited from the vision encoder.
Table~\ref{tab:probes} reports functional probe $R^2$ for key representation sources.
Two complementary effects characterize LLM processing.

\paragraph{Distributed degradation.}
Averaging across image tokens, the LLM degrades both geometric and functional action information relative to the projector.
RSA drops from $0.615$ to $0.501$ ($\Delta = 0.114$); probe $R^2$ drops from $0.706$ to $0.683$ (95\% CI of $\Delta$: $[0.011, 0.037]$, $p = 0.001$).
Figure~\ref{fig:layer_curves} shows the per-layer trajectory: the sharpest drop occurs at the embedding interface (Projector $\to$ Layer~1), with non-monotonic fluctuations through middle layers before partial recovery near the output layers.

\paragraph{Readout concentration.}
While distributed image tokens lose information, the last pre-action token (\texttt{last\_preaction}) concentrates action-relevant information beyond the projector baseline.
Trained readout achieves $R^2 = 0.731$, exceeding SigLIP+Projector ($R^2 = 0.706$, $p = 0.001$; Figure~\ref{fig:probe_bars}).
This concentration is training-dependent: untrained readout ($R^2 = 0.686$) falls below the projector baseline, while trained readout ($0.731$) substantially exceeds it ($\Delta = 0.045$, $p < 0.0001$).

Together, these results characterize the LLM as a \emph{redistribution} mechanism: training teaches the model to route action-relevant information from distributed image tokens to the readout position, sacrificing distributed fidelity to concentrate a functionally superior action representation at the decoding interface~\cite{haon2025}.

\begin{table}[t]
\centering
\caption{Functional probe $R^2$ ($\uparrow$) for action prediction from frozen representations. The LLM degrades distributed representations (image\_mean) relative to the projector but concentrates superior action information at the readout token---an effect that requires training. DINOv2 matches the projector in $R^2$ despite much lower RSA.}
\label{tab:probes}
\small
\begin{tabular}{lcc}
\toprule
Representation & $R^2$ & $p$ vs.\ Proj. \\
\midrule
SigLIP+Projector & 0.706 & --- \\
VLA Trained image\_mean (L19) & 0.683 & 0.001 \\
VLA Trained readout (L19) & \textbf{0.731} & 0.001 \\
VLA Untrained readout (L32) & 0.686 & --- \\
DINOv2 cls\_token & 0.712 & 0.220 (n.s.) \\
\bottomrule
\end{tabular}
\end{table}

\begin{figure}[t]
\centering
\includegraphics[width=\columnwidth]{figures/fig3.pdf}
\caption{Per-layer RSA ($\rho$) with the action space across all 32 LLM layers under four experimental conditions (trained/untrained $\times$ with/without instruction). The sharpest drop occurs at the embedding interface (Proj $\to$ L1). Middle layers show non-monotonic fluctuations. Dashed lines: SigLIP+Projector ($\rho{=}0.615$) and DINOv2 ($\rho{=}0.388$) reference levels.}
\label{fig:layer_curves}
\end{figure}

\begin{figure}[t]
\centering
\includegraphics[width=\columnwidth]{figures/fig4.pdf}
\caption{Functional probe $R^2$ for action prediction. The LLM degrades distributed image\_mean ($R^2{=}0.683$) below the projector baseline ($0.706$, $p{=}0.001$) but concentrates superior information at the trained readout ($0.731$, $p{=}0.001$). Untrained readout ($0.686$) falls below the projector, confirming that concentration requires training ($\Delta{=}0.045$, $p{<}0.0001$). DINOv2 ($0.712$) is statistically indistinguishable from the projector ($p{=}0.22$).}
\label{fig:probe_bars}
\end{figure}

%------------------------------------------------------------------
\subsection{Training Routes, Does Not Create, Action Geometry}
\label{sec:training_routing}

This experiment tests whether LLM training creates new action-relevant geometric structure or merely routes pre-existing structure inherited from the vision encoder.
Table~\ref{tab:training_effects} reports the effect of training across metrics.
The $2 \times 2$ factorial analysis (Figure~\ref{fig:factorial}) reveals a striking dissociation: training has no detectable effect on geometric structure but a large effect on functional concentration.

\paragraph{Geometric invariance to training.}
Training does not alter the LLM's geometric alignment with the action space.
For image\_mean RSA, the training effect is $\Delta = +0.021$ (95\% CI $[-0.031, +0.064]$, $p = 0.20$, n.s.); for last\_preaction RSA, $\Delta = +0.008$ ($[-0.040, +0.056]$, $p = 0.37$, n.s.).
This null result is consistent with the vision encoder---whose weights are identical between conditions ($\cos\text{sim} = 1.0$)---being the sole source of action-relevant geometry.

\paragraph{Geometric invariance to instruction.}
Language instructions similarly fail to modulate geometric structure.
Across the $2 \times 2$ matrix, the instruction effect is $\Delta = 0.004$ (n.s.).
A dedicated control experiment with four instruction conditions---normal task instructions ($\rho = 0.489$), random instructions from other tasks ($0.484$), word-shuffled instructions ($0.483$), and empty strings ($0.472$)---confirms content invariance: all values fall within $\pm 0.017$ of the mean (Figure~\ref{fig:factorial}b).

\paragraph{Functional concentration requires training.}
In contrast to the geometric null results, readout probe $R^2$ shows a large training effect: $\Delta = +0.045$ (CI $[+0.032, +0.058]$, $p < 0.0001$).
Training transforms the readout from below-projector ($R^2 = 0.686$) to above-projector ($R^2 = 0.731$) without altering the underlying geometry.

These results support a \emph{routing} interpretation of LLM training in VLAs~\cite{brohan2023_rt2}: the LLM learns to concentrate pre-existing vision-encoder geometry at the readout position for action decoding, rather than constructing new action-relevant representations.

\begin{table}[t]
\centering
\caption{Training and instruction effects on action-relevant structure. Training does not alter geometric alignment (RSA, both n.s.)\ but enables functional concentration at the readout token ($R^2$, $p < 0.0001$). The projector-to-image\_mean $R^2$ drop confirms distributed degradation.}
\label{tab:training_effects}
\small
\begin{tabular}{lccc}
\toprule
Comparison & $\Delta$ & 95\% CI & $p$ \\
\midrule
Training (image\_mean RSA) & $+$0.021 & [$-$0.031, $+$0.064] & 0.20 \\
Training (readout RSA) & $+$0.008 & [$-$0.040, $+$0.056] & 0.37 \\
Training (readout $R^2$) & $+$0.045 & [$+$0.032, $+$0.058] & $<$0.0001 \\
Proj.\ vs.\ image\_mean ($R^2$) & $+$0.023 & [$+$0.011, $+$0.037] & 0.001 \\
\bottomrule
\end{tabular}
\end{table}

\begin{figure}[t]
\centering
\includegraphics[width=\columnwidth]{figures/fig5.pdf}
\caption{\textbf{(a)}~Training $\times$ Instruction matrix for image\_mean RSA ($\rho$). Neither training ($\Delta{=}0.021$, $p{=}0.20$) nor instruction ($\Delta{=}0.004$, n.s.)\ produces a statistically detectable effect. \textbf{(b)}~Instruction content control: normal ($0.489$), random ($0.484$), shuffled ($0.483$), and empty ($0.472$) instructions all yield RSA within $\pm 0.017$ of the mean $\mu{=}0.482$, confirming that instruction content does not modulate distributed geometric structure.}
\label{fig:factorial}
\end{figure}

\begin{figure}[t]
\centering
\includegraphics[width=\columnwidth]{figures/fig6.pdf}
\caption{Comparison of three similarity metrics across representation sources. \textbf{(a)}~RSA (Spearman $\rho$) and \textbf{(c)}~RDM kernel CKA produce consistent rankings (Proj $>$ Trained $>$ Untrained $>$ DINOv2). \textbf{(b)}~Linear CKA~\cite{kornblith2019} partially reverses the ordering (Untrained $>$ Trained $>$ Proj), illustrating that metric choice affects conclusions about representation similarity.}
\label{fig:metric_comparison}
\end{figure}
